\documentclass{article}
\usepackage{xcolor}
\usepackage[utf8]{inputenc}
\usepackage{amsmath}
\usepackage{graphicx}
\graphicspath{{Imagenes/}}

\title{\huge \color{red}Taller de Herramientas Computacionales}
\author{Brenda Paola García Rivas}
\date{24.01.19}

\begin{document}
	\maketitle
	\begin{center}
		\includegraphics [scale=1.5]{1.png}
	\end{center}

	\newpage
	\title{\huge Jueves 24 de Enero. Clase 14}\\
	
	Hoy fue una clase muy interesante, se inició con el problema de fibonacci, el cuál es un problema de recursividad; ayer se vió como resolverlo, sin embargo, esta vez el profesor nos explicó una forma más fácil de hacer los llamados para cada número. 
	Este método implica guardar los resultados en una lista, en este se modificará que: el programa observe que esté un número en la lista para después guardarlo y posteriormente usarlo, pero ¿cómo se sabe que el décimo, onceavo, o el número que sea está en la lista?, pues fácil, teniendo en cuenta la longitud de la lista.
	
	Además de esto, también se resolvieron algunas dudas de la tarea, se inició con el problema de la creación de un laberinto, para realizar éste, debemos tener en mente varias cosas:
	\begin{enumerate}
		\item 
	Primero, se identificó en el pizarrón cómo se iba a hacer el laberinto, se intentó hacerlo lo más facil posible dado que sería únicamente un ejmplo. El laberinto tenía forma de 3x3, en dónde únicamente se podía pasar horizantalmente por la fila de enmedio.
		\item 
	Ahora bien, lo siguiente en cuestionarse es: ¿Puedo avanzar hacia enfrente?, si la respuesta es "sí", entonces, se debe mover una columna, y en la coordenada 'y' se le aumentaría un lugar.
		\item 
	Lo siguiente que hay que ver es que 
	\end{enumerate}
	\begin{verbatim}
		# -*- coding: utf-8 -*-
		L=[[True, True, True, True],    #Se define el laberinto
		[False, False, False, True],
		[True, True, False, True]]
		def resolver(L,e):
		print (e)
		n=len(L[0])  #Columna 1
		x=e[0]
		y=e[1]
		if y==n-1: salida
		return e[0]+1,e[1]+1  #Ya llegué
		else: #Si no se ha llegado...
		if L[x][y+1] == False: #Se puede mover una columna más, por eso se aumenta a y 
		e=[x,y+1]
		return resolver(L,e)
		else:
		print 'Ya no se puede avanzar'
	\end{verbatim}
\end{document}